\section{Exactly uniform demands}
\label{sec:reduction}

We finish by converting the metric obstruction into capacities with unit
demand on every unordered pair.  We first record the precise duality
statement, allowing distinct points to have distance zero. This is the
average-distortion form of the cut/embedding correspondence
\cite[Section~4]{LLR95}\cite[Section~1.3]{Rab08}; the explicit programs
below keep the unit-demand normalization visible.

\begin{lemma}[Cut-cone duality]
\label{lem:duality}
Let $d$ be a semimetric on $[n]$, where $n\geq 2$, that satisfies the
triangle inequalities and is a squared Hilbert distance.  Suppose that
$a,b>0$ satisfy
\begin{equation}
\label{red:averages}
\frac{1}{n^2}\sum_{i,j=1}^n d(i,j)\geq a,
\qquad
\frac{1}{n^2}\sum_{i,j=1}^n
\|F(i)-F(j)\|_1\leq b
\end{equation}
for every map $F$ into a finite-dimensional $\ell_1$ space such that
$\|F(i)-F(j)\|_1\leq d(i,j)$ for all $i,j$.
Then there are nonnegative capacities $c_{ij}$ on unordered pairs of
distinct vertices such that, with demand exactly one on every pair,
\begin{equation}
\label{red:gap-lemma}
\OPT(C)\geq 1,
\qquad
0<\GL(C)\leq \frac{b}{a}.
\end{equation}
\end{lemma}

\begin{proof}
Let $\mathcal S$ be the finite collection of nonempty proper subsets of
$[n]$.  For $S\in\mathcal S$, write
$\delta_S(i,j)=|\mathbf 1_S(i)-\mathbf 1_S(j)|$.  Consider the finite
linear program.\footnote{Matou\v{s}ek and
Rabinovich~\cite{MR01} study the related convex hull of individually
dominated line metrics. Here the cut coefficients are arbitrary
nonnegative numbers; only their sum as a distance function is required
to be dominated by $d$.}
\begin{equation}
\label{red:primal}
\begin{aligned}
\text{maximize}\quad &
\sum_{S\in\mathcal S}|S|(n-|S|)t_S,\\
\text{subject to}\quad &
\sum_{S\in\mathcal S}t_S\delta_S(i,j)\leq d(i,j)
\quad (i<j),\\
&t_S\geq 0\quad(S\in\mathcal S).
\end{aligned}
\end{equation}
It is feasible by taking every $t_S=0$.  For any feasible $t$, the map
$F_t(i)=(t_S\mathbf 1_S(i))_{S\in\mathcal S}$ takes values in
$\ell_1^{|\mathcal S|}$ and is a contraction for $d$.  Its sum of
distances over unordered pairs is exactly the objective in
\eqref{red:primal}.  Consequently \eqref{red:averages} bounds the primal
optimum by $n^2b/2$.  The primal optimum is attained: for each
$S\in\mathcal S$ and any pair separated by $S$, feasibility implies
$0\leq t_S\leq d(i,j)$, so its feasible region is closed and bounded.

The dual of \eqref{red:primal}, with one nonnegative variable $c_{ij}$
for each unordered pair, is
\begin{equation}
\label{red:dual}
\begin{aligned}
\text{minimize}\quad &\sum_{i<j}c_{ij}d(i,j),\\
\text{subject to}\quad &
\sum_{i\in S,\ j\notin S}c_{ij}\geq |S|(n-|S|)
\quad(S\in\mathcal S),\\
&c_{ij}\geq 0\quad(i<j).
\end{aligned}
\end{equation}
Here the cut sum counts each separated unordered pair once.
The dual is feasible, for example with $c_{ij}=1$ for every pair.
Finite-dimensional linear programming duality therefore gives an
optimal dual solution $C$ satisfying
\begin{equation}
\label{red:dual-objective}
\sum_{i<j}c_{ij}d(i,j)\leq \frac{n^2b}{2}.
\end{equation}
Its cut constraints give $\OPT(C)\geq1$ with exactly unit pair demands.

Set $D=\sum_{i<j}d(i,j)\geq n^2a/2>0$.  If
$d(i,j)=|z_i-z_j|^2$ in a Hilbert space, the vectors
$z_i/\sqrt D$ give the feasible normalized SDP distances $d(i,j)/D$.
The span of the finitely many $z_i$ has dimension at most $n$, so these
vectors can be realized in $\R^n$.  The triangles are preserved by
normalization, and \eqref{red:dual-objective} gives
$\GL(C)\leq b/a$.

Finally, the graph of pairs with $c_{ij}>0$ is connected: otherwise the
vertex set of one component would violate a cut constraint in
\eqref{red:dual}.  Choose a spanning tree $\mathcal T$ in this graph
and put $\gamma=\min_{\{i,j\}\in\mathcal T}c_{ij}>0$.  For any feasible
SDP distance array $d'$, the triangle inequalities along tree paths
imply
\[
1=\sum_{i<j}d'(i,j)
\leq \binom n2\sum_{\{i,j\}\in\mathcal T}d'(i,j).
\]
Hence its objective is at least $\gamma/\binom n2>0$.  This proves
$\GL(C)>0$, including when the original semimetric $d$ has zero
distances between distinct vertices.
\end{proof}

\subsection{Replacing the parameter distribution by multiplicities}

Fix a sufficiently large integer $m$, and use the vertex set $V$,
semimetric $d$, and chart map $v_1$ from \Cref{prop:metric}.  Define the
probability weights
\[
w_v=\Pr_{\theta\ \mathrm{uniform\ on}\ T}
\big[v_1(\theta)=v\big],
\qquad T=(-1,1)^m.
\]
Weights on vertices outside the first chart are thus zero.
By \Cref{prop:metric,prop:contraction}, for absolute positive constants
that may change from line to line,
\begin{equation}
\label{red:weighted-bounds}
\begin{gathered}
|V|\leq \exp(Cm\log m),\qquad
\operatorname{diam}(V,d)\leq Cm,\qquad
\sum_{v,z\in V}w_vw_zd(v,z)\geq cm,\\
\sum_{v,z\in V}w_vw_z\|F(v)-F(z)\|_1
\leq C\sqrt m\,(\log m)^{5/2}
\end{gathered}
\end{equation}
for every finite-dimensional $\ell_1$-valued contraction $F$ for $d$.

Every chart vertex must remain present, even if its weight is zero:
its distance constraints still restrict contractions on the first chart.
We therefore give each vertex at least one copy. In the choice below,
the term $m^2|V|$ controls the measure error, while $m^m$ ensures the
lower bound on $\log n$ needed in the final parameter conversion. Set
\[
\begin{gathered}
M=m^m+m^2|V|,
\qquad k_v=1+\lfloor Mw_v\rfloor,\\
\widetilde V=\{(v,h):v\in V,\ 1\leq h\leq k_v\},
\qquad n=|\widetilde V|.
\end{gathered}
\]
Give these copies the semimetric
$\widetilde d((v,h),(z,k))=d(v,z)$.
It is again a squared Hilbert distance satisfying all triangles, by
assigning the same Hilbert vector to every copy of a vertex.  In
particular, distances between copies of one vertex are zero.

Let $\mu_v=k_v/n$ be the pushforward to $V$ of the uniform probability
measure on $\widetilde V$.  To bound the rounding error, write
$e_v=k_v-Mw_v\in(0,1]$ and $E_0=\sum_v e_v$.  Then
$n=M+E_0$, where $0<E_0\leq |V|$, and
\[
\mu_v-w_v=\frac{e_v-E_0w_v}{n}.
\]
Using the convention
$\|\mu-w\|_{\mathrm{TV}}=\frac12\sum_v|\mu_v-w_v|$, we obtain
\begin{equation}
\label{red:tv}
\|\mu-w\|_{\mathrm{TV}}
\leq \frac{E_0}{n}
\leq \frac{|V|}{M}
\leq \frac{1}{m^2},
\qquad
\|\mu\otimes\mu-w\otimes w\|_{\mathrm{TV}}
\leq \frac{2}{m^2}.
\end{equation}
The second inequality follows by writing the difference of the product
measures as
$(\mu-w)\otimes\mu+w\otimes(\mu-w)$ and applying the triangle
inequality for total variation.

For any function with values in $[0,Cm]$, the change in its expectation
under these two product measures is at most $2C/m$.
Applying this to $d$ gives
\begin{equation}
\label{red:uniform-metric}
\frac{1}{n^2}\sum_{x,y\in\widetilde V}\widetilde d(x,y)
=\sum_{v,z\in V}\mu_v\mu_zd(v,z)
\geq c'm
\end{equation}
for all sufficiently large $m$.
Every contraction $\widetilde F$ from $(\widetilde V,\widetilde d)$
into a finite-dimensional $\ell_1$ space agrees on copies of a vertex,
because the distance between those copies is zero.  Since $k_v\geq1$
for every $v$, it therefore induces a contraction $F$ on all of $V$.
The function $\|F(v)-F(z)\|_1$ is bounded by $d(v,z)\leq Cm$.
Thus \eqref{red:weighted-bounds} and \eqref{red:tv} give
\begin{equation}
\label{red:uniform-contraction}
\frac{1}{n^2}\sum_{x,y\in\widetilde V}
\|\widetilde F(x)-\widetilde F(y)\|_1
\leq C'\sqrt m\,(\log m)^{5/2}.
\end{equation}

Apply \Cref{lem:duality} to
\eqref{red:uniform-metric}--\eqref{red:uniform-contraction}.
After labeling $\widetilde V$ by $[n]$, this gives nonnegative
capacities $C^{(m)}$ with unit demands on all distinct pairs and
\begin{equation}
\label{red:m-gap}
\frac{\OPT(C^{(m)})}{\GL(C^{(m)})}
\geq c''\frac{\sqrt m}{(\log m)^{5/2}},
\qquad \GL(C^{(m)})>0.
\end{equation}

\subsection{The number of vertices}

Since $M\leq n\leq M+|V|$, \eqref{red:weighted-bounds} implies, for
an absolute constant $K$ and all sufficiently large $m$,
\begin{equation}
\label{red:size}
m\log m\leq\log n\leq K m\log m.
\end{equation}
In particular $n\to\infty$ and $\log\log n\geq\log m$ for large $m$.
Consequently
\[
\frac{\sqrt{\log n}}{(\log\log n)^3}
\leq \sqrt K\,
\frac{\sqrt m}{(\log m)^{5/2}}.
\]
Combining this with \eqref{red:m-gap} proves \Cref{thm:main}.
